\section{Problem 4: Gear Train Design}

\subsection{Problem description}
The objective of this exercise is to choose a gear train design that achieves a specific transmission ratio.
We define the gear train with four gears with the numbers of teeth $N_1, N_2, N_3, N_4$.
The target transmission ratio is: 
\begin{equation}
R = \frac{N_1 \cdot N_3}{N_2 \cdot N_4}
\end{equation}

Here, the target transmission ratio is: $R_{target} = 1/6.931 \approx 0.1442793$.
Each gear must have an integer number of teeth between 12 and 60.
The problem is difficult because the search space is discrete and nonlinear, meaning classical optimization cannot be used directly.

Therefore, two heuristic optimization algorithms are tested: 
\begin{enumerate}
    \item Particle Swarm Optimization (PSO)
    \item Simulated Annealing (SA)
\end{enumerate}

The goal is to compare their performance and their behavior on this discrete optimization problem.

\subsection{Cost function}
The cost function measures the error between the achieved ratio and the target ratio:
\begin{equation}
f(N_1,N_2,N_3,N_4 ) = \left| \frac{N_1 \times N_3}{N_2 \times N_4} - R_{target} \right|
\end{equation}
A good solution is one where this error approaches zero.

\subsection{Results}

\textbf{Particle Swarm Optimisation (PSO):}
\begin{itemize}
    \item Best combination: $[N_1, N_2, N_3, N_4] = [18, 60, 25, 52]$
    \item Achieved Ratio: 0.1442308
    \item Absolute Error: $4.856 \times 10^{-5}$
\end{itemize}

\textbf{Simulated Annealing (SA):}
\begin{itemize}
    \item Best combination: $[37, 21, 30, 17]$
    \item Achieved Ratio: 3.109
    \item Absolute Error: 2.965
\end{itemize}

\subsection{Work discussion}
The Particle Swarm Optimisation successfully found a high-quality solution with a very low error ($< 10^{-4}$). The discrete nature of the problem (integer teeth) was handled by rounding the particle positions.

However, the Simulated Annealing run failed to converge to a valid solution close to the target in this instance, ending with a large error. This suggests that the cooling schedule (cooling rate 0.9) or initial temperature (1000) might need adjustment, or the neighbor generation function (random perturbation) requires tuning to better explore the discrete integer space of the gear teeth. PSO proved more robust for this specific discrete problem configuration.
